\chapter{Introduction}
\label{ch:introduction}

\section{Purpose and Scope}

The SIRC-1 CPU is sold on its own, so it is documented in three manuals that stand apart from any system built around
it. SIRCADA is made for the SIRCUS system and has a manual of its own. Together the four manuals have the following
scopes:

\begin{table}[H]
	\centering
	\begin{tabular}{lp{8.2cm}}
		\toprule
		\textbf{Manual} & \textbf{Scope}                                                                                                          \\
		\midrule
		\archmanual{}   & ISA, instruction encodings, canonical assembly notation, exceptions, programmer-visible state                           \\
		\asmmanual{}    & Complete assembly syntax, directives, symbols, meta-instructions, and toolchain usage                                   \\
		\hwmanual{}     & Pinout, bus interface, timing, electrical characteristics                                                               \\
		\videomanual{}  & The picture processing unit: package and pinout, memory map, registers, data formats, video modes, timing (this manual) \\
		\bottomrule
	\end{tabular}
	\caption{The SIRC-1 and SIRCUS Manual Set}
	\label{tab:manual-set}
\end{table}

This manual covers SIRCADA as a component: its package and pins, what the CPU can read and write, how the device turns
the contents of its video memory into a picture, and the timing of its signals. It does not describe the CPU, its
instructions or its bus protocol. The device is a bus peripheral like any other, so the \hwmanual{} is authoritative
for how the CPU makes a bus cycle, and this manual is authoritative for how the device answers it and for what it does
once it has been reached.

\section{Device Overview}

The video device draws a picture one scanline at a time, with no frame buffer: for every group of four pixels it
fetches tilemap entries and tile data from its video memory, resolves them through a palette, and sends the result to
the display.

\begin{table}[H]
	\centering
	\begin{tabular}{lp{9cm}}
		\toprule
		\textbf{Property} & \textbf{Value}                                                                                                    \\
		\midrule
		Picture size      & 256 by 224 visible pixels (262 lines NTSC, 312 PAL; Chapter~\ref{ch:timing})                                      \\
		Colour            & 15 bits per pixel, 5 bits for each of red, green and blue                                                         \\
		Palette           & 256 entries of 16 bits, directly addressable by the CPU                                                           \\
		Video memory      & 32,768 words (64 KB), shared by tile data and tilemaps                                                            \\
		Tile data         & 4 bits per pixel                                                                                                  \\
		Backgrounds       & Four orthogonal in Mode 1; two isometric and one orthogonal in Mode 2; two orthogonal and one isometric in Mode 3 \\
		Sprites           & 128 entries, three layers, 32 per line (Chapter~\ref{ch:sprites})                                                 \\
		Interrupts        & Two: vertical sync (NMI) and line compare (IRQ4)                                                                  \\
		Line timing       & 1364 master clocks per line, 262 lines per frame                                                                  \\
		\bottomrule
	\end{tabular}
	\caption{Video Device Summary}
	\label{tab:device-summary}
\end{table}

\section{Video Modes}
\label{sec:video-modes}

The device has three video modes, selected by the \texttt{MODE} field of the Base Configuration register
(Section~\ref{sec:reg-base-config}). The modes share the memory map, the registers, the palette and the data formats.
They differ in how the backgrounds are drawn and in the fetch schedule that supports them.

\begin{table}[H]
	\centering
	\begin{tabular}{llp{7.5cm}}
		\toprule
		\textbf{Mode} & \textbf{\texttt{MODE} value} & \textbf{Description}                                                                                                 \\
		\midrule
		Mode 1        & 0                            & Four orthogonal backgrounds of square tiles and sprites. The reset mode.                                             \\
		Mode 2        & 1                            & Two isometric backgrounds: square tiles displayed as diamonds on a 2:1 ground plane, and an orthogonal status layer. \\
		Mode 3        & 2                            & One isometric background behind two orthogonal backgrounds.                                                          \\
		\bottomrule
	\end{tabular}
	\caption{Video Modes}
	\label{tab:video-modes}
\end{table}

Mode 1 is the conventional mode and is described in Chapter~\ref{ch:mode1}. Mode 2 is a specialised mode for games set
on a flat map seen from above at an angle. It is good at that and at little else; Chapter~\ref{ch:mode2} describes what
it does and, as importantly, what it gives up. Mode 3 (Chapter~\ref{ch:mode3}) keeps one isometric background and gives
the fetch slots of the other to two orthogonal ones. The values 3--7 of \texttt{MODE} are reserved; software must not
write them. The \texttt{MODE} field is sampled at the first cycle of each line, so a change takes effect at the start
of the next line.

\section{Block Diagrams}
\label{sec:block-diagrams}

Figure~\ref{fig:block-system} shows the device as the CPU sees it: four areas of memory behind one bus interface, a
video output, and two interrupt lines. Figure~\ref{fig:block-ppu} shows what is inside. The chapters that follow
describe each block; the figure says how they feed each other.

\begin{figure}[H]
	\centering
	\begin{tikzpicture}[
			scale=0.88, transform shape,
			block/.style={draw, rounded corners=2pt, fill=sirclightgray, align=center, font=\small},
			mem/.style={draw, fill=white, align=center, font=\footnotesize, minimum height=0.9cm, minimum width=5.4cm},
			flow/.style={-{Stealth[length=2mm]}, thick},
			note/.style={font=\scriptsize, text=sircgray}
		]
		\node[block, minimum width=2.2cm, minimum height=3.8cm] (cpu) at (0,0) {SIRC-1\\CPU};
		\node[block, minimum width=3.0cm, minimum height=3.8cm] (bus) at (3.7,0) {Bus interface\\and address\\decode\\(segment\\\texttt{0x000C})};
		\node[mem] (m1) at (9.3,2.1) {Registers \texttt{0x0000}--\texttt{0x001A}};
		\node[mem] (m2) at (9.3,0.7) {Palette \texttt{0x6000}, 256 words};
		\node[mem] (m3) at (9.3,-0.7) {Sprite table \texttt{0x7000}, 256 words};
		\node[mem] (m4) at (9.3,-2.1) {Video memory \texttt{0x8000}, 32,768 words};
		\node[block, minimum width=2.8cm, minimum height=4.4cm] (ppu) at (14.2,0) {Picture\\processing\\unit};
		\node[block, minimum width=2.8cm] (out) at (14.2,-3.9) {Video output\\256 by 224};
		\draw[flow, <->] (cpu) -- node[above, note]{address, data} (bus);
		\draw[thick] (5.6,2.1) -- (5.6,-2.1);
		\draw[thick] (bus.east) -- (5.6,0);
		\foreach \y in {2.1,0.7,-0.7,-2.1} {
				\draw[flow, <->] (5.6,\y) -- (6.6,\y);
				\draw[flow, <->] (12.0,\y) -- (ppu.west |- 0,\y);
			}
		\draw[flow] (ppu) -- (out);
		\draw[flow] (ppu.north) -- ++(0,0.9) -| node[above, near start, note]{NMI (vertical sync), IRQ4 (line compare)} (cpu.north);
	\end{tikzpicture}
	\caption{The Video Device in the System}
	\label{fig:block-system}
\end{figure}

\begin{figure}[H]
	\centering
	\begin{tikzpicture}[
			scale=0.86, transform shape,
			block/.style={draw, rounded corners=2pt, fill=sirclightgray, align=center, font=\footnotesize, minimum height=0.95cm, minimum width=3.1cm},
			store/.style={draw, fill=white, align=center, font=\footnotesize, minimum height=0.95cm, minimum width=3.1cm},
			flow/.style={-{Stealth[length=2mm]}, thick},
			ctl/.style={-{Stealth[length=2mm]}, thick, dashed},
			note/.style={font=\scriptsize, text=sircgray}
		]
		% Register file runs down the left edge: every block reads it
		\node[block, minimum width=8.6cm, minimum height=1.0cm, rotate=90] (regs) at (-1.5,3.4) {Register file (read by every block)};
		% Control row
		\node[block] (intr) at (1.6,8.0) {Interrupt logic\\NMI, IRQ4};
		\node[block] (timing) at (5.3,8.0) {Timing generator\\line, group, slot};
		\node[block] (seq) at (9.3,8.0) {Fetch sequencer\\slot scheduler};
		% Column 1: memories and CPU
		\node[store] (vram) at (1.6,5.4) {Video memory};
		\node[store] (cpu) at (1.6,3.2) {CPU access\\(vertical blanking only)};
		\node[store] (oam) at (5.3,-0.9) {Sprite table};
		% Column 2: arbiter
		\node[block, minimum height=5.0cm] (arb) at (5.3,3.0) {Memory\\arbiter\\\scriptsize renderer first};
		% Column 3: engines
		\node[block] (flat) at (9.3,5.4) {Orthogonal layer engines\\(Modes 1, 3, status)};
		\node[block] (iso) at (9.3,3.8) {Isometric layer engines\\(Modes 2, 3)};
		\node[block] (que) at (9.3,2.2) {Status layer queue\\(Mode 2)};
		\node[block] (spe) at (9.3,0.6) {Sprite engine\\scan, fetch, line buffer};
		% Column 4: compositor and output
		\node[block, minimum height=5.0cm, minimum width=2.9cm] (comp) at (13.3,3.0) {Compositor\\priority\\S1 B1 S2 B2\\S3 B3 B4\\shading};
		\node[store] (pal) at (13.3,-0.9) {Palette\\(colour lookup)};
		\node[block] (bri) at (13.3,-2.5) {Brightness\\and blanking};
		\node[block] (vo) at (13.3,-4.1) {Video output};
		% Control flows
		\draw[ctl] (timing) -- (intr);
		\draw[ctl] (timing) -- (seq);
		\draw[ctl] (seq.south) -- ++(0,-0.55) -| node[above, near start, note]{slot grants} (arb.north);
		% Data flows
		\draw[flow, <->] (vram) -- (vram.east -| arb.west);
		\draw[flow] (cpu) -- (cpu.east -| arb.west);
		\draw[flow] (oam.east) -| ([xshift=-1.1cm]spe.south);
		\foreach \e in {flat,iso,que} { \draw[flow] (arb.east |- \e) -- (\e.west); }
		\draw[flow] (arb.east |- spe) -- (spe.west);
		\foreach \e in {flat,iso,que,spe} { \draw[flow] (\e.east) -- (\e.east -| comp.west); }
		\draw[flow] (comp.south) -- (pal.north);
		\draw[flow] (pal) -- (bri);
		\draw[flow] (bri) -- (vo);
	\end{tikzpicture}
	\caption{Blocks of the Picture Processing Unit}
	\label{fig:block-ppu}
\end{figure}

\section{Conventions}

\begin{itemize}
	\item All addresses and sizes are in 16-bit words, as on the CPU. A word address is an offset from the base of the device's
	      segment unless stated otherwise.
	\item Bit 0 is the least significant bit. Bit diagrams draw the most significant bit on the left.
	\item Hexadecimal values carry a \texttt{0x} prefix and upper case digits.
	\item Ranges of values and addresses use an en dash: \texttt{0x0000}--\texttt{0x00FF}.
	\item \textbf{Reserved} fields and registers are set aside by the architecture. Software must not rely on what they
	      read, and should write zero to a reserved field it does not otherwise use.
	\item A value that the design has not fixed is \textbf{not yet specified}. It is not undefined, and no value is implied.
\end{itemize}
